\section{When signed force changes the survival test}
\label{sec:diagnostics}

\subsection{An unsigned comparison}

The signed estimate enters the survival criterion through the
integral \(E_\beta(\Delta)\). On the same box, the bound
\(\lvert g_x\rvert\le K\) already gives \(B\le K\). If a signed bound
\(\beta_{\mathrm s}\) is available, one can use
\(\beta=\min\{K,\beta_{\mathrm s}\}\).
For \(\Delta>0\), \(E_\beta(\Delta)=\int_0^\Delta e^{4\beta t}\,dt\)
is strictly increasing in \(\beta\). Thus a signed bound can improve
the gap test on a fixed box only when \(\beta_{\mathrm s}<K\). The
following implication identifies boxes where the unsigned comparison
already ensures a positive gap margin.

\begin{proposition}[Redundancy of the signed gap test]
\label{thm:redundancy}
Suppose the starting data and uniform deflated-factor bounds of the
survival theorem hold. Assume the two strict confinement inequalities in
\eqref{eq:survival-margins}. If
\begin{equation}
 Q_-Q_+>32D^2\Delta,
 \label{eq:redundancy-screen}
\end{equation}
the gap condition in \eqref{eq:survival-margins} holds with \(\beta=K\).
\end{proposition}

\begin{proof}
Since \(K\ge0\), upper-gap confinement gives
\[
 E_K(\Delta)\le\Delta e^{4K\Delta}
       <\frac{4D^2\Delta}{Q_+}.
\]
Hence \(Q_--8E_K(\Delta)>Q_- - 32D^2\Delta/Q_+>0\).
The integral bound includes \(K=0\).
\end{proof}

The converse fails in the class of real polynomial heat solutions.
For
\[
 h(\tau,z)=(z-X)^2+2(\tau-T),\qquad T=\tfrac14,
\]
the exact deflated factor is \(1\), so \(A=K=B=0\),
\(c=X\) and \(q(\tau)=8(T-\tau)\) on the real-pair side.
Choose \(Q_-=Q_+=2\), \(D=1\), \(\Delta=3/16\),
\(R_0=0\), and \(X>R+1\) with \(R>0\). The two confinement
inequalities hold, and \(Q_--8E_0(\Delta)=1/2>0\), although
\(Q_-Q_+=4<6=32D^2\Delta\).
At the exact collision horizon \(\Delta=T\), the sharp initial
bound instead gives \(Q_--8E_0(T)=0\). For any \(D^2>2T\)
ensuring strict initial upper-gap confinement,
\(Q_-Q_+=64T^2<32D^2T\), so the redundancy condition
also fails at that horizon. For this implementation a failed screen
therefore leaves the direct unsigned gap test to be evaluated.
The scale \(q(0)/8\) in the quadratic is a reference scale for a
tight box with \(4D^2\) comparable to \(q(0)\).

\subsection{Conditioning of separate uniform constants}

Even with a small source error, the numerator and denominator
entering \(A\) in \eqref{eq:force-constants} may be constrained at
different endpoints. At an actual simple pair let \(\delta=b-a>0\)
and write \(h=(z-a)(z-b)U(z)\), where \(U\) here is the unscaled
deflated factor. Differentiation gives
\begin{equation}
 U(a)=-\frac{h_x(a)}{\delta},\qquad
 U(b)=\frac{h_x(b)}{\delta},\qquad
 U_x(a)=\frac{2U(a)-h_{xx}(a)}{2\delta}.
 \label{eq:conditioning-jets}
\end{equation}
The same box bounds both endpoints, so \(M_1+e_1\ge\lvert U_x(a)\rvert\)
and \(0<m-e_0\le\lvert U(b)\rvert\), with the common sign chosen
as in the transfer lemma. It follows that
\begin{equation}
 A\ge \frac{\lvert U_x(a)\rvert}{\lvert U(b)\rvert},
 \qquad K\ge A^2.
 \label{eq:uniform-cost}
\end{equation}
Strict numerical inequalities in opposite directions for the two
endpoint quantities can therefore force a lower bound on the
\emph{allowances} \(A\) and \(K\), independent of the truncation
order. Their quotient in \eqref{eq:uniform-cost} is invariant under
common scalar normalization. It measures the cost of using separate
uniform constants. A pointwise relative estimate can preserve the
correlation between the numerator and its own denominator.

For comparison, the actual root jets give the identity
\[
 c'=-\frac12\left(\frac{h_{xx}(a)}{h_x(a)}
                 +\frac{h_{xx}(b)}{h_x(b)}\right)
\]
on a simple-root slab. Bounds for these quotients require paid slope
floors, and can produce a smaller midpoint allowance than \(2A\) on
a slab where those floors are available. An exterior energy budget can
also be compared in the same normalization. In Newman time,
\(q_\lambda=8+4qB\) by \eqref{eq:gap-law}. For the ordered two-root
energy \(\mathcal E=2/q\) this yields
\begin{equation}
 \mathcal E_\lambda=-4\mathcal E^2+B_{\mathrm{RT}},
 \qquad B_{\mathrm{RT}}=-\frac{8B}{q}.
 \label{eq:rt-conversion}
\end{equation}
Equations \eqref{eq:conditioning-jets}--\eqref{eq:rt-conversion}
allow the comparison to use the same root pair, interval, and force
normalization.

\subsection{The direction of the Lehmer sum}

For this comparison, assume at a fixed slice that the exterior
product gives
\[
 g_x(x)=-\sum_{\rho\ne a,b}(x-\rho)^{-2}
\]
on the closed segment of the actual simple pair. In a finite product
this follows by logarithmic differentiation. For a genus-one product
the linear exponential factors contribute zero to \(g_x\), and the
squared reciprocal tail ensures absolute and uniform convergence on
the segment. Indeed, for sufficiently distant zeros
\(\inf_{x\in[a,b]}|x-\rho|\ge|\rho|/2\), and the finitely many nearer
exterior zeros are separated from the segment. The sign argument
requires every exterior zero to be real and outside \([a,b]\).
These extra hypotheses concern this comparison alone.

Set \(\delta=b-a\), \(q=\delta^2\), and
\[
 G=\sum_{\rho\ne a,b}
      \bigl((a-\rho)^{-2}+(b-\rho)^{-2}\bigr),
 \qquad
 S=\sum_{\rho\ne a,b}\frac{1}{(a-\rho)^2(b-\rho)^2}.
\]
The mean of \(-(x-\rho)^{-2}\) over \([a,b]\) is
\(-1/((a-\rho)(b-\rho))\). Put \(u=a-\rho\) and \(v=b-\rho\).
Then
\[
 u^{-2}+v^{-2}-2(uv)^{-1}
     =\frac{(v-u)^2}{u^2v^2}.
\]
Summing gives
\begin{equation}
 G+2B=qS\ge0,\qquad -G/2\le B\le0.
 \label{eq:lehmer-identity}
\end{equation}
The upper bound is strict at slices with at least one real exterior
zero, since each product \(uv\) is positive. It supplies no
quantitative negative margin by itself.

The information direction is fixed by \eqref{eq:gap-law}.
An upper bound \(G\le G_*\) gives \(B\ge-G_*/2\) and hence
\(q'\le-8+2G_*q\). A survival lower bound uses \(B\le\beta\)
to obtain \(q'\ge-8-4\beta q\). A value of \(G\) at one slice
also leaves its behavior over the rest of the slab open. To
improve the upper bound \(B\le0\) within the real-product class,
one would need additional information, for instance a lower
bound for a positive contribution to \(-B\) from known exterior
zeros. Real-axis source estimates can also give an upper force
bound without the real-product hypothesis.

With complex conjugate exterior zeros the identity can persist
whenever the product formula and convergence still hold. Conjugation
makes the combined contribution real, with no fixed sign. For a conjugate pair
\(\rho,\bar\rho\), write \(P=(a-\rho)(b-\rho)\). Its contributions
to \(S\) and \(B\) are \(2\operatorname{Re}(P^{-2})\) and
\(-2\operatorname{Re}(P^{-1})\), respectively. For example, the
real-polynomial snapshot with \(a=1,b=3\) and
\(U(z)=(z-1)^2+1\) has \(P=-1-2i\) at \(\rho=1+i\).
It gives \(S=-6/25\) and \(B=2/5\), exhibiting failure of
both signs in \eqref{eq:lehmer-identity} in the complex-exterior
class.

\subsection{A quartic collision with negative force}

Consider the forward heat evolution
\begin{equation}
 h(\tau,z)=z^4+(12\tau-5)z^2+12\tau^2-10\tau+4 .
 \label{eq:lehmer-quartic}
\end{equation}
Both \(h_\tau\) and \(h_{zz}\) equal \(12z^2+24\tau-10\).
At zero time the roots are \(-2,-1,1,2\). With \(a=1,b=2\),
the distances to the two exterior zeros are \((2,3)\) and \((3,4)\).
Consequently
\[
 q=1,\qquad G=\frac14+\frac29+\frac1{16}=\frac{77}{144},
 \qquad B=-\frac16-\frac1{12}=-\frac14,\qquad S=\frac5{144}.
\]
Thus \(qG<4/5\), the algebraic Lehmer diagnostic, and
\(G+2B=qS>0\). The positive pair nonetheless collides. To see this,
the quadratic in \(w=z^2\) has discriminant
\[
 \mathcal D(\tau)=96\tau^2-80\tau+9,\qquad
 T=\frac{10-\sqrt{46}}{24}
\]
with \(T\) its first zero. For \(0\le\tau<T\) the two \(w\)-roots
are real and distinct. Their sum is \(5-12\tau>0\) and their
product is \(12(\tau-5/12)^2+23/12>0\), so both are positive.
At \(T\) the polynomial is \((z^2-r^2)^2\), where
\(r^2=\sqrt{46}/4\). In particular \(h_{zz}(T,r)=2\sqrt{46}>0\).
The event is a nondegenerate local double at \(r\), with a separate
reflected double at \(-r\). The discriminant is negative just
after \(T\).

Throughout \(0\le\tau<T\), the negative pair contributes
\[
 B=-\frac{1}{2a(a+b)}-\frac{1}{2b(a+b)}=-\frac{1}{2ab}<0 .
\]
The law with zero exterior force and \(q(0)=1\) would close at
\(1/8\). Here \(q'>-8\) before the event and \(1/8<T<1/4\).
Thus the favorable force delays this collision while the gap
still closes. The backward-heat field
\(F_\lambda(z)=h(1/4-\lambda,z)\) has its collision at
\(\lambda_*=(\sqrt{46}-4)/24>0\). This control concerns
polynomial heat solutions. Its collision rules out simultaneously
meeting the strict survival hypotheses on a box containing the
pair up to the event.

\subsection{Normalization of the historical data}

For the heat integral in Stopple's notation \cite{Stopple}, the source
is \(2\Phi(u/2)\), where \(\Phi\) is the kernel used here. With
\(\widehat H_t\) denoting Stopple's deformation, the substitution
\(u=2w\) gives
\[
 \widehat H_t(\gamma)
   =\int_0^\infty e^{tu^2}2\Phi(u/2)\cos(\gamma u)\,du
   =4 H_{4t}(2\gamma).
\]
Since \(\Xi=2\widehat H_0\), this also yields
\(H_0(x)=\Xi(x/2)/8\). Thus \(x=2\gamma\) and \(\lambda=4t\).
The endpoint sum in ordinate units is \(G_\gamma=4G_x\),
while \(q_\gamma=q_x/4\), so the diagnostic \(qG\) is invariant.
